% 【本文件写什么】api-v0 契约层，叶子，只写语义不写实现
% - crate 路径：os/components/wateros-fs/fs-devfs/devfs-api/api-v0/src/
% - 列出并说明：pub trait、核心类型、错误枚举、常量
% - 每个公开 API 的前置条件、后置条件、线程/中断上下文限制
% - 与 Linux/用户态约定的对齐点与 deliberate 差异
% - v0 版本当前行为与后续可替换 extension 点
% 【事实来源】
% - 上述 crate 下全部 .rs 源文件
% - docs/exports/public-api/ 中对应符号表，以源码为准
% 【不写什么】
% - impl 内部数据结构、汇编、具体算法步骤

\subsubsection{api-v0}

\paragraph{类型。}
\code{DevNodeType}区分\code{Block}、\code{Character}、\code{Unsupported}。\code{DevNode}含逻辑路径，例如 \code{/dev/vda}，与类型，供启动日志与块设备路径解析。错误复用\code{fs\_api\_v0::FsError}/\code{FsResult}，与根卷挂载错误映射一致。

\paragraph{\code{DevFsManager}契约。}
\begin{itemize}
 \item \code{refresh()}：按当前平台驱动枚举重建节点表与块/字符绑定；调用后\code{list\_nodes()}反映最新快照；
 \item \code{set\_dt\_unsupported\_paths(paths)}：登记DTB解析出但尚无驱动实现的占位路径，类型为\code{Unsupported}；
 \item \code{list\_nodes()}：返回节点列表，顺序由实现决定；
 \item \code{register\_block\_device}/\code{lookup\_block\_device}：路径与\code{SharedBlockDevice}双向绑定；
 \item \code{register\_character\_device}/\code{lookup\_character\_device}：字符设备绑定；
 \item \code{default\_root\_block\_path()}：默认用于根卷probe的块设备路径；无设备时\code{None}。
\end{itemize}

\paragraph{与\code{FsImpl}的关系。}
devfs管理面与块卷\code{FsImpl}解耦：\code{fs-impl/impl-devfs}仅将devfs登记为\code{FsKind::Other("devfs")}能力项，\code{probe}恒为\code{None}，不参与块设备格式识别。真实根卷probe由ext4等块FS impl完成。
